\documentclass{article}

\title{Пример документа LaTeX}
\author{Автор}
\date{\today}

\begin{document}

\maketitle

\section{Введение}
LaTeX -- это мощная система компьютерной вёрстки, которая особенно популярна в научных кругах для подготовки документов, содержащих математические формулы и библиографию. В отличие от обычных текстовых редакторов, LaTeX фокусируется на структуре документа, позволяя автору сосредоточиться на содержании, а не на его внешнем виде.

В этом документе мы рассмотрим основные элементы типичной статьи, включая введение, основную часть и заключение. LaTeX предоставляет широкие возможности для форматирования текста, вставки формул, таблиц, изображений и других элементов.

\section{Основная часть}
В основной части документа обычно размещается основное содержание. Это может быть описание методов исследования, анализ данных, теоретические выкладки или другие материалы, зависящие от темы статьи.

LaTeX позволяет легко создавать списки, как нумерованные:
\begin{enumerate}
\item Первый пункт
\item Второй пункт
\item Третий пункт
\end{enumerate}

Так и маркированные:
\begin{itemize}
\item Пункт 1
\item Пункт 2
\end{itemize}

\section{Заключение}
В заключении подводятся итоги работы. LaTeX остается незаменимым инструментом для подготовки научных публикаций, диссертаций, курсовых работ и других документов, требующих сложного форматирования. Его использование позволяет создавать документы высокого качества с последовательным и профессиональным оформлением.

При правильном использовании LaTeX значительно упрощает процесс подготовки документов, особенно тех, которые содержат большое количество формул, таблиц и библиографических ссылок.

\end{document}